\section{Formal Mathematics}

This appendix gives the theory a compact formal core. It states, in precise symbols, what the substrate is, what
a fill is, how a beat acts, and what is conserved, and then records the principal emergent results as formal
propositions. The body of the paper is the intuition and the evidence. This is the skeleton.

\subsection{Notation, clear symbols}

We avoid single Greek letters as variables, since they are easy to confuse. We use plain Latin letters and short
words. The only Greek kept is the numeric constant $\pi = 3.14159\ldots$, which is universal. The dictionary below
also lists the conventional Greek name used elsewhere in the literature, in brackets, so the reader can connect
the two.

\begin{table}[h]
\centering
\begin{tabular}{lll}
\toprule
symbol & meaning & (usual name) \\
\midrule
$M$ & the vibe mesh, the $\{3,4,3,4\}$ honeycomb & \\
$V$ & the set of cells of the mesh & \\
$v_1, v_2$ & cells, the vibes, elements of $V$ & \\
$C$ & the $24$ directions of a cell, the coin & \\
$c$ & a direction, an element of $C$ & \\
$n(v, c)$ & the neighbour of $v$ across face $c$ & \\
$T = \{-1, 0, +1\}$ & the tone alphabet, the three tone values & \\
$t(v, c)$ & the tone in slot $c$ of cell $v$, in $T$ & \\
$f$ & a fill, a full configuration, a map $V \times C \to T$ & \\
$\mathsf{Stream}$ & the global transport map on fills & \\
$\mathsf{Collide}$ & the local interaction map on fills & \\
$R = \mathsf{Collide} \circ \mathsf{Stream}$ & the rule, one beat & \\
$B$ & the beat, the discrete tick & \\
$F = \{M, B, T, R\}$ & the flow, the four base things together & \\
$b$ & the beat index, in $\mathbb{Z}$ & \\
$Q$ & total charge, the sum of all tones & \\
$\mathsf{dist}(v_1, v_2)$ & graph distance, the step count & \\
$r(n)$ & the ripple level, the number of cells at distance $n$ from a root & \\
$\mathsf{gr}$ & the ripple growth rate & \\
$H$ & the expansion rate, $\ln(\mathsf{gr})/3$ & (Hubble) \\
$\mathsf{cc}$ & the cosmological constant, $3 H^2$ & ($\Lambda$) \\
$E, k, m$ & energy, wavenumber, mass of a mode & \\
$\mathsf{wk}$ & the weak mixing fraction, $\sin^2\theta_W$ & ($\theta_W$) \\
$\mathsf{decay}$ & a correlator decay exponent & ($\alpha$) \\
\bottomrule
\end{tabular}
\caption{The clear-symbol dictionary for this appendix.}
\end{table}

\subsection{The substrate}

\begin{definition}[The flow]
The base of the theory is the \emph{flow}, the tuple $F = \{M, B, T, R\}$ of the mesh $M$, the beat $B$, the tone
alphabet $T$, and the rule $R$. The four are specified in turn below. Everything else in the theory is a property
of $F$ or an emergent consequence of iterating the rule $R$ beat by beat.
\end{definition}

\begin{definition}[The mesh]
The \emph{mesh} $M$ is the regular hyperbolic honeycomb $\{3,4,3,4\}$ tiling $\mathbb{H}^4$. Its cells form a set
$V$. Each cell is a $24$-cell, a four-dimensional regular solid with $24$ octahedral faces, so each cell shares
each of its $24$ faces with exactly one other cell. The \emph{directions} $C$ are the $24$ face-normals of a cell,
identified with the $24$ roots of the root system $D_4$. We write $|C| = 24$.
\end{definition}

\begin{definition}[The neighbour map and distance]
There is a fixed map $n : V \times C \to V$ sending $(v, c)$ to the cell sharing face $c$ with $v$. It is an
involution on the face it crosses, $n(n(v, c), \, \overline{c}) = v$ for the opposite direction $\overline{c}$.
The mesh graph has vertex set $V$ and an edge $\{v, n(v, c)\}$ for each $v, c$. Its shortest-path metric is
$\mathsf{dist}$. Distance is the only notion of space at the base, space is step-counting on this graph.
\end{definition}

\begin{definition}[Growth]
Fix a root cell $o$. The ripple levels are $r(n) = |\{v : \mathsf{dist}(o, v) = n\}|$. The mesh is
\emph{generated}, not given, the cell at the root exists first and each beat adds the next ripple, so the universe
at beat $b$ is the ball of radius $b$. This identifies radial distance with cosmic time, one ripple per beat.
\end{definition}

\subsection{The fill}

\begin{definition}[The fill]
A \emph{fill} is a map $f : V \times C \to T$. The value $f(v, c)$, also written $t(v, c)$, is the tone in
directional slot $c$ of cell $v$, read as the charge currently leaving $v$ through face $c$. The set of all fills
is $T^{\,V \times C}$. A single cell has local fill space $T^{C}$ of size $3^{24}$. There is no other state, in
particular no value stored on edges, the relation across a face is computed from the two slots and never stored.
\end{definition}

\begin{definition}[Scalar tone and charge]
The \emph{scalar tone} of cell $v$ in fill $f$ is $\mathsf{tone}(v) = \sum_{c \in C} t(v, c)$, an integer in
$[-24, 24]$. The \emph{total charge} is
\[
Q(f) \;=\; \sum_{v \in V} \; \sum_{c \in C} t(v, c),
\]
the sum of all slots over the whole mesh, an integer.
\end{definition}

\subsection{The dynamics}

A beat is two maps composed, first transport, then interaction.

\begin{definition}[Stream]
For each cell $v$ and direction $c$, streaming moves the slot value along its own line of travel into the
neighbour. Writing $\overline{c}$ for the opposite direction, the streamed fill is
\[
(\mathsf{Stream}\, f)(v, c) \;=\; f\big(n(v, \overline{c}), \, c\big),
\]
that is, the value now in slot $c$ of $v$ is the value that the back-neighbour was sending forward along line $c$.
Each slot has exactly one source and one destination, so $\mathsf{Stream}$ is a bijection of fills, a global
permutation of slots. It carries no decision, a charge keeps its direction, which is momentum.
\end{definition}

\begin{definition}[Collide]
There is a fixed local map $\mathsf{col} : T^{C} \to T^{C}$ on the $24$ slots of one cell. It is applied at every
cell simultaneously,
\[
(\mathsf{Collide}\, f)(v, \cdot) \;=\; \mathsf{col}\big(f(v, \cdot)\big).
\]
The local map $\mathsf{col}$ is required to be (i) a \emph{bijection} of $T^{C}$, hence reversible, and (ii)
\emph{charge-conserving} on the cell, $\sum_c \mathsf{col}(w)_c = \sum_c w_c$ for every local vector $w$. It is
built from four move-types, pass-through, annihilate, create, and scatter. Its two-slot shadow is the reversible,
charge-conserving permutation of the nine pair-states $(-1,-1), (+1,+1), (-1,0), (0,-1), (+1,0), (0,+1), (0,0),
(+1,-1), (-1,+1)$, with the create-flip-annihilate three-cycle $(0,0) \to (+1,-1) \to (-1,+1) \to (0,0)$ and two
inert fixed points.
\end{definition}

\begin{definition}[The rule and a beat]
One \emph{beat} is the composition
\[
R \;=\; \mathsf{Collide} \circ \mathsf{Stream} : T^{\,V \times C} \to T^{\,V \times C}.
\]
The fill at beat $b$ from initial fill $f_0$ is $f_b = R^{\,b}(f_0)$. Time is the iteration index $b$, there is no
external clock.
\end{definition}

\subsection{Conservation and reversibility}

\begin{proposition}[Conservation of charge]
For every fill $f$, $\; Q(R\, f) = Q(f)$. Hence $Q$ is constant along every trajectory.
\end{proposition}

\textit{Proof.} $\mathsf{Stream}$ permutes the slot values without changing their multiset, so it preserves the
global sum $Q$. $\mathsf{Collide}$ preserves the per-cell sum by requirement (ii), and the global sum is the sum
of per-cell sums, so it too preserves $Q$. The composition preserves $Q$. Checked to machine precision in
\cite{pollard2026vibetest}. \hfill$\square$

\begin{proposition}[Reversibility]
$R$ is a bijection of fills, with inverse
$R^{-1} = \mathsf{Stream}^{-1} \circ \mathsf{Collide}^{-1}$. Hence every fill has a unique past and a
unique future, and $R^{-N}(R^{N}(f)) = f$ exactly.
\end{proposition}

\textit{Proof.} $\mathsf{Stream}$ is a bijection by construction, and $\mathsf{Collide}$ is a bijection because
the local map $\mathsf{col}$ is, applied independently per cell. The composition of bijections is a bijection.
Verified forward-then-back to zero error at scale in \cite{pollard2026vibetest}. \hfill$\square$

\begin{result}[The arrow]
Reversibility makes create and annihilate inverse moves on one cycle. The \emph{arrow} is the preference for the
forward sense of that cycle, the create direction $(0,0) \to (+1,-1)$, which turns a balanced peace state into a
charge pair. Because $Q$ is conserved, every created pair is net zero, a $+1$ with a $-1$. The arrow is therefore
the only ingredient that selects a temporal orientation, and it adds no charge.
\end{result}

\subsection{The coin algebra}

\begin{definition}[The coin and its decomposition]
The $24$ directions $C$ are the roots of $D_4$. Under the cell-symmetry group they decompose as
\[
24 \;=\; 8_v \,\oplus\, 8_s \,\oplus\, 8_c,
\]
the vector representation of $\mathrm{SO}(8)$ and its two spinor representations, permuted by the order-$6$
triality symmetry $S_3$.
\end{definition}

\begin{proposition}[Spinor requirement forces $\{3,4,3,4\}$]
A collision built on $C$ carries a half-integer-spin (spinor) representation if and only if the direction set
contains a spinor summand. The $\{3,4,3,4\}$ coin does, by the $8_s, 8_c$ above. The $12$ directions of the $3$D
candidate $\{5,3,4\}$ decompose as $1 \oplus 3 \oplus 3' \oplus 5$, all integer spin, with no spinor. Hence among
the candidate substrates, $\{3,4,3,4\}$ alone supports spin at the base.
\end{proposition}

\begin{result}[Odd-charge solitons are fermions]
A matter soliton is a hopfion, a smooth tone-texture with a Hopf map $S^3 \to S^2$ of integer linking degree. A
soliton of odd degree acquires the sign $-1$ under a full $2\pi$ rotation, by the Finkelstein-Rubinstein and
Wilczek-Zee arguments, so it obeys Fermi statistics. The spin-statistics link is therefore geometric, not
assumed. Triple-checked by linking degree, configuration-space fundamental group, and direct exchange
\cite{pollard2026vibetest}.
\end{result}

\subsection{Emergent propagation}

The continuous, complex, quantum description is the coarse-grained limit of the discrete real walk, never a base
ingredient. The following are measured dispersions of the rule.

\begin{result}[Light cone]
A pulse front advances exactly one cell per beat, a finite frame-independent maximum speed, normalized to
$1$. In the massless case the rule restricted to one line is the pure shift, and the measured dispersion is
\[
E(k) \;=\; k \quad \text{(exact, no continuity assumed)}.
\]
\end{result}

\begin{result}[Massive Dirac dispersion]
Adding a scatter (a coin rotation by mass angle $m$) gives the measured lattice Dirac dispersion
\[
\cos E \;=\; \cos m \,\cos k, \qquad E(0) = m,
\]
so rest energy equals the mass, and mass is the coarse-grained mixing rate. In the small-$k$, small-$m$ limit
this is the relativistic relation $E^2 = k^2 + m^2$.
\end{result}

\subsection{Geometry, growth, and cosmology}

\begin{proposition}[Ripple recurrence and growth rate]
The ripple levels $r(n)$ satisfy a fixed linear recurrence whose characteristic polynomial is that of the
splitting matrix of the $\{3,4,3,4\}$ self-similar tree. Its dominant (Perron) root is the growth rate
$\mathsf{gr}$. Measured ripples $1, 24, 456, 8376, \ldots$ give an asymptotic ratio $\mathsf{gr} \approx 18.3$,
while an early-level average on a truncated ball gives $\mathsf{gr} \approx 11$.
\end{proposition}

\begin{result}[de Sitter expansion]
With physical-space volume growing as the cube of a scale factor $a$, and the cell count growing by $\mathsf{gr}$
per beat,
\[
a(b) \;\sim\; \mathsf{gr}^{\,b/3} \;=\; e^{H b}, \qquad H = \frac{\ln \mathsf{gr}}{3},
\]
so the expansion is exponential, $a''/a = H^2 > 0$, with cosmological constant $\mathsf{cc} = 3 H^2$. Using the
early-level $\mathsf{gr} \approx 11$ gives $H \approx 0.80$ and $\mathsf{cc} \approx 1.9$ in beat units. The flat
twin $\{3,4,3,3\}$ grows polynomially with ratio about $2.5$ and has no exponential expansion, so $H$ is a
curvature effect.
\end{result}

\begin{definition}[The cusp]
A \emph{horosphere} is a level set of a Busemann function toward an ideal boundary point of $\mathbb{H}^4$. It is
intrinsically flat. The ideal $\{4,3,4\}$ cubic element of $\{3,4,3,4\}$ is such a horosphere, the \emph{cusp},
combinatorially the Euclidean cubic lattice $\mathbb{Z}^3$. Physical space is the cusp. Its measured spectral
dimension is $2.97$, against $2.16$ for a generic rough slice.
\end{definition}

\subsection{Gauge and the Standard-Model numbers}

\begin{proposition}[$\mathrm{SO}(10)$ from coin and tone]
Adjoining the ternary tone sign as one extra weight axis to the $D_4 = \mathrm{SO}(8)$ root system yields the
$D_5 = \mathrm{SO}(10)$ root system. The $16$-dimensional spinor of $\mathrm{SO}(10)$ is exactly one fermion
generation, split as $16 = 8_s$ (tone $+$) $\oplus\, 8_c$ (tone $-$). Triality gives three generations.
\end{proposition}

\begin{proposition}[Automatic breaking to the Standard Model]
Any self-condensate (a Higgs in the $16$) leaves unbroken the subgroup $\mathrm{SU}(5)$ (twenty unbroken roots),
and the Georgi-Glashow chain $\mathrm{SO}(10) \to \mathrm{SU}(5) \to \mathrm{SU}(3) \times \mathrm{SU}(2) \times
\mathrm{U}(1)$ follows, with $16 = 1 \oplus 10 \oplus \overline{5}$, the complete matter content of one
generation.
\end{proposition}

\begin{result}[The weak mixing fraction]
Summing the weak isospin and charge over one generation,
\[
\mathsf{wk} \;=\; \frac{\sum T_3^2}{\sum q^2} \;=\; \frac{3}{8} \quad \text{exactly, at the unification scale,}
\]
where $T_3$ is the third weak-isospin component and $q$ the electric charge over the sixteen states. Running this
down with standard renormalization reaches the measured $0.231$ for a supersymmetric low-energy content.
\end{result}

\begin{result}[Mass relations]
At the unification scale, the bottom and tau masses are equal, $m_b = m_\tau$, and the determinant relation
$\det(M_{\mathrm{lepton}}) = \det(M_{\mathrm{down}})$ holds, the latter from the discrete trace identity
$\mathrm{Tr}\, Y = 0$ over the sixteen, where $Y$ is the hypercharge. These are ratios, not absolute masses.
\end{result}

\subsection{Gravity, as an equation of state}

\begin{result}[Newtonian potential on the cusp]
The static lattice Green's function $G(r)$ of the cusp Laplacian, computed by the difference method
$G(r) - G(r_0)$ to cancel the zero-mode contamination, fits
\[
G(r) - G(2) \;=\; \frac{a_1}{r} + a_0, \qquad a_1 \approx 0.083 \approx \frac{1}{4\pi},
\]
a clean Newtonian $1/r$ in the $3$D cusp.
\end{result}

\begin{result}[Holographic boundary correlator]
On the infinite tree of branching number $q$, the boundary-to-boundary correlator computed by the Bethe cavity
recursion decays as $1/r^{\mathsf{decay}}$ with
\[
\mathsf{decay} \;=\; \frac{2 \ln(1/w)}{\ln q}, \qquad w = \frac{1}{q} \;\Rightarrow\; \mathsf{decay} = 2,
\]
a clean $1/r^2$, confirmed independently by a tree-tensor contraction at $1/r^{1.93}$.
\end{result}

\begin{result}[The Einstein field equations]
Take the area-law entropy $\mathsf{Ent} = \mathsf{Area}/4$ (the holographic $1/4$) and the local horizon
temperature $\mathsf{Temp} = \kappa / (2\pi)$ from the Unruh response, where $\kappa$ is the surface gravity.
Requiring the Clausius balance $\delta\mathsf{Heat} = \mathsf{Temp}\, \delta\mathsf{Ent}$ across every local
horizon, for all null directions, is equivalent to the full field equations
\[
\mathsf{Ein}(x) + \mathsf{cc}\,\mathsf{g}(x) \;=\; 8\pi\, \mathsf{Stress}(x),
\]
where $\mathsf{Ein}$ is the Einstein tensor, $\mathsf{g}$ the metric, and $\mathsf{Stress}$ the matter
stress-energy. The coefficient $8\pi$ is forced by the area-law $1/4$ and is the same Newton constant as in the
$1/r$ law. The weak field reduces to Poisson's equation $\nabla^2 \mathsf{pot} = 4\pi\, \mathsf{density}$, and a
photon deflects by $4 G M / x$, twice the naive Newtonian value.
\end{result}

\begin{result}[Black-hole thermodynamics]
With $\mathsf{Ent} = \mathsf{Area}/4 = 4\pi M^2$ and $\mathsf{Temp} = 1/(8\pi M)$ for a mass-$M$ horizon, the
first law $\mathrm{d}M = \mathsf{Temp}\,\mathrm{d}\mathsf{Ent}$ holds, the Smarr identity $M = 2\,\mathsf{Temp}\,
\mathsf{Ent}$ holds, the Bekenstein bound is saturated, the heat capacity $\mathrm{d}M/\mathrm{d}\mathsf{Temp} =
-8\pi M^2$ is negative, and the evaporation lifetime scales as $M^3$. The substrate's own de Sitter horizon obeys
the same laws with $\mathsf{Temp} = H/(2\pi)$ and $\mathsf{Ent} = \pi/H^2$.
\end{result}

\subsection{Electromagnetism}

\begin{result}[Discrete Gauss law and the gauge field]
Per-cell charge conservation, charge in equals charge out each beat, is a discrete Gauss law, the conserved
current sourcing an emergent $\mathrm{U}(1)$. A vortex link field gives a single-plaquette flux $\mathsf{flux}$
whose Wilson loop (the sum of link phases around the plaquette) equals $\mathsf{flux}$ and is invariant under a
gauge shift $A \to A + \mathrm{d}\theta$, and a charge encircling it acquires exactly the phase $\mathsf{flux}$
(Aharonov-Bohm), while a non-enclosing loop acquires $0$. The non-abelian case replaces phases by $\mathrm{SO}(3)$
link matrices with non-commuting holonomy.
\end{result}

\subsection{Addressing and computation}

\begin{proposition}[Self-describing addresses]
Fix a root $o$ and the breadth-first spanning tree of the mesh graph. Each cell $v$ has a unique address, the
sequence of child-indices on the path $o \to v$, of length $\mathsf{dist}(o, v) = O(\log |\,\text{ball}\,|)$.
Crucially the tree has \emph{no same-level edges}, every neighbour of a cell is a parent or a child, so the parent
is dropping the last digit and the children are appending an admissible digit. The only non-tree edges are
confluences, a cell reachable as a child of more than one parent, and they form a finite-state transducer on a
bounded window of the generator address (deterministic at window two over the verified ripples).
\end{proposition}

\begin{result}[Universality]
The substrate is computationally universal. (i) The signed-majority form of the collision, with a $+1$ bias and
two $-1$ fills, computes $\mathrm{NAND}$, which is functionally complete, so it builds any Boolean circuit and in
particular the universal elementary rule $110$. (ii) A reversible, charge-conserving register machine
(increment, decrement-or-jump, halt) runs on the mesh with registers held as charge in addressed subtrees, where
increment is the create move and decrement is annihilation, computing correctly with $Q$ conserved. (iii) The
cusp $\mathbb{Z}^3$ runs Conway's Life, a strongly universal cellular automaton. So the substrate is weakly
universal by (i) and (ii), and strongly universal by (iii).
\end{result}

\subsection{The discipline, and the recorded negatives}

\begin{principle}[Nothing added]
Every result above is a property of the four base things of the flow $F = \{M, B, T, R\}$ or an emergent
consequence of them. No fifth ingredient is used. Where a phenomenon does not emerge, the negative is recorded as
a formal statement, not hidden.
\end{principle}

\begin{result}[Recorded negatives]
The following are recorded negatives, established by the same standard as the positives.
\begin{enumerate}
\item \emph{No spontaneous selves from generic initial data.} The pure rule from a generic fill churns, a
localized seed disperses (radius of gyration grows without bound, recentered pattern identity falls to zero), so
durable selves do not crystallize from equilibrium initial data without seeding or topological protection.
\item \emph{No genuine self-tower.} The coarse-grained net-charge persistence rises under coarse-graining on both
the cusp and the bulk radial tree, but it does \emph{not} beat a plain-diffusion control, so it is a generic
conserved-field slow mode, not hierarchical selfhood. The correct self-measure (identity through material
turnover) confirms the cusp supports propagating self-patterns while the pure rule does not make one.
\item \emph{No scale-free avalanches.} The ballistic light cone forbids branching-process self-organized
criticality, so hierarchical integration, where it occurs, comes from fast ballistic routing, not sandpile
cascades.
\end{enumerate}
\end{result}

The one named candidate for a possible fifth ingredient is higher-order binding, the clean binding of structure
\emph{between} units rather than within one. It sits behind the open holes, cross-gap attraction, super-selves,
and the nesting tower. It is named, not added, and not yet proven either to emerge from the four or to be
irreducible.
